\section*{附加习题: 拓扑性质与 $\pi_{1}$}
\label{ch13-extra}

上一节的结论说明: 对于 $B$ 的覆叠空间的分类而言, 空间 $B$ 是道路连通、局部道路连通且半局部单连通的空间的假定是适当的. 此处我们将说明: 对于 $B$ 的各种拓扑性质与 $B$ 的基本群之间的关系而言, 对空间 $B$ 的以上假定也是合适的.

\begin{exercise}
\label{exercise:ch13-extra-01}
设 $X$ 是一个空间, $\mathcal{A}$ 是 $X$ 的一个开覆盖. 在何种条件下, 存在 $X$ 的一个开覆盖 $\mathcal{B}$ 加细 $\mathcal{A}$, 并且满足条件: 对于 $\mathcal{B}$ 中每一对相交非空的元素 $B$, $B'$, 有 $B\cup B'$ 包含于 $\mathcal{A}$ 的某一元素之中?
\begin{enumerate}[(a)]
\item 证明: 若 $X$ 是一个可度量化空间, 那么满足上述条件的开覆盖 $\mathcal{B}$ 是存在的. [提示: 选取 $\varepsilon(x)$ 使得 $B(x,\ 3\varepsilon(x))$ 包含于 $\mathcal{A}$ 的某一元素之中. 设 $\mathcal{B}$ 由所有开集 $B(x,\ \varepsilon(x))$ 组成.]
\item 证明: 若 $X$ 是一个紧致的 Hausdorff 空间, 那么满足上述条件的开覆盖 $\mathcal{B}$ 是存在的. [提示: 设 $A_{1}$, $\cdots$, $A_{n}$ 是覆盖 $X$ 的 $\mathcal{A}$ 的一个有限子族. 选取 $X$ 的一个开覆盖 $C_{1}$, $\cdots$, $C_{n}$ 使得对于每一个 $i$ 有 $\overline{C}_{i}\subset A_{i}$. 对于 $\{1,\ \cdots,\ n\}$ 的每一个非空子集 $J$, 考虑集合
\[
B_{J}=\bigcap_{j\in J}A_{j}-\bigcup_{j\notin J}\overline{C}_{j}.]
\]
\end{enumerate}
\end{exercise}

\begin{exercise}
\label{exercise:ch13-extra-02}
\begin{theorem}[第二可数基下基本群可数]
设空间 $X$ 是道路连通、局部道路连通且半局部单连通的. 如果 $X$ 是具有可数基的正则空间, 则 $\pi_{1}(X,\ x_{0})$ 是可数的.
\end{theorem}
证明: 设 $\mathcal{A}$ 是 $X$ 的由道路连通的开集 $A$ 所构成的一个覆盖, 使得对于每一个 $A\in\mathcal{A}$ 及每一个 $a\in A$, 由内射所诱导的同态 $\pi_{1}(A,\ a)\to\pi_{1}(X,\ a)$ 是平凡的. 设 $\mathcal{B}$ 是满足习题 1 中所述条件的 $X$ 的由非空道路连通集合所构成的一个可数开覆盖. 对于每一个 $B\in\mathcal{B}$ 选取一点 $p(B)\in B$. 对于 $\mathcal{B}$ 中的每一对满足 $B\cap B'\neq\varnothing$ 的元素 $B$, $B'$, 选取 $B\cup B'$ 中从 $p(B)$ 到 $p(B')$ 的一条道路 $g(B,\ B')$. 我们称 $g(B,\ B')$ 为\textbf{选择道路}.

设 $B_{0}$ 为 $\mathcal{B}$ 中的一个固定元素, $x_{0}=p(B_{0})$. 若 $f$ 是 $X$ 中的一条以 $x_{0}$ 为基点的回路, 那么依以下步骤证明 $f$ 道路同伦于一些选择道路的乘积:
\begin{enumerate}[(a)]
\item 证明存在 $[0,\ 1]$ 的一个分划
\[
0=t_{0}<\cdots<t_{n}=1
\]
使得 $f$ 将 $[t_{n-1},\ t_{n}]$ 映入 $B_{0}$, 并且对于每一个 $i=1$, $\cdots$, $n-1$, $f$ 将 $[t_{i-1},\ t_{i}]$ 映入 $\mathcal{B}$ 的某一个元素 $B_{i}$ 中. 设 $B_{n}=B_{0}$.
\item 设 $f_{i}$ 为从 $[0,\ 1]$ 到 $[t_{i-1},\ t_{i}]$ 的正线性映射与 $f$ 的复合, $g_{i}=g(B_{i-1},\ B_{i})$. 选取 $B_{i}$ 中从 $f(t_{i})$ 到 $p(B_{i})$ 的一条道路 $\alpha_{i}$, 若 $i=0$ 或 $i=n$, 设 $\alpha_{i}$ 为 $x_{0}$ 处的常值道路. 证明
\[
[f_{i}]*[\alpha_{i}]=[\alpha_{i-1}]*[g_{i}].
\]
\item 证明 $[f]=[g_{1}]*\cdots*[g_{n}]$.
\end{enumerate}
\end{exercise}

\begin{exercise}
\label{exercise:ch13-extra-03}
设 $p\colon E\to X$ 是一个覆叠映射, 使得 $\pi_{1}(X,\ x_{0})$ 是可数的. 证明: 若 $X$ 是具有可数基的正则空间, 则 $E$ 也是具有可数基的正则空间. [提示: 设 $\mathcal{B}$ 是由 $X$ 的道路连通集所构成的一个可数基. 对于 $B\in\mathcal{B}$, 设 $\mathcal{C}$ 是由 $p^{-1}(B)$ 的所有道路连通分支组成的族. 参见第 53 节中的习题 6.]
\end{exercise}

\begin{exercise}
\label{exercise:ch13-extra-04}
\begin{theorem}[紧致空间的基本群有限生成]
设空间 $X$ 是道路连通、局部道路连通且半局部单连通的. 如果 $X$ 是紧致的 Hausdorff 空间, 则 $\pi_{1}(X,\ x_{0})$ 是有限生成的, 从而也是可数的.
\end{theorem}
证明: 重复习题 2 中所述的证明. 选取 $\mathcal{B}$ 是一个有限族. 与先前一样, 我们有等式
\[
[f]=[g_{1}]*\cdots*[g_{n}].
\]
对于每一个 $x\in X$, 选取一条从 $x_{0}$ 到 $x_{1}$ 的道路 $\beta_{x}$, 设 $\beta_{x_{0}}$ 为常值道路. 若 $g=g(B,\ B')$, 定义
\[
L(g)=\beta_{x}*(g*\overline{\beta}_{y}),
\]
其中 $x=p(B)$ 并且 $y=p(B')$. 证明:
\[
[f]=[L(g_{1})]*\cdots*[L(g_{n})].
\]
\end{exercise}

\begin{exercise}
\label{exercise:ch13-extra-05}
设 $X$ 是无限耳环 (参见第 80 节中的例 1). 证明 $X$ 是具有可数基的紧致的 Hausdorff 空间, 其基本群含有不可数多个元素. [提示: 设 $r_{n}\colon X\to C_{n}$ 是一个收缩. 给定 $0$ 和 $1$ 的一个序列 $a_{1}$, $a_{2}$, $\cdots$, 证明存在 $X$ 中的一条回路 $f$, 使得对于每一个 $n$, 元素 $(r_{n})_{*}[f]$ 是平凡的当且仅当 $a_{n}=0$.]
\end{exercise}
